\documentclass[10pt,a4paper]{amsart}
\usepackage[margin=19mm]{geometry}
\usepackage{amsmath,amssymb,mathtools}
\usepackage[scr=rsfs,cal=euler]{mathalfa}
\usepackage{xfrac}
\usepackage[unicode,hidelinks,bookmarks=false]{hyperref}
\newcommand{\ie}{i.e.\ }
\newcommand{\cf}{cf.\ }
\newcommand{\resp}{respectively\ }
\newcommand{\Subsection}{\subsection}
\providecommand{\nobreakdash}{\nobreak\hbox{-}}
\setlength{\emergencystretch}{2em}
\multlinegap0pt
\usepackage{amssymb}
\usepackage{tikz}
\usepackage{bm}
\usepackage{enumerate}
\usepackage[shortlabels]{enumitem}
\usepackage{csquotes}
\usepackage{xcolor}
\newtheorem{prop}{Proposition}
\newcommand*{\abs}[1]{\lvert #1\rvert}
\title{On the extrema of the mean subtree order of graphs}
\author{Stijn Cambie, Jorik Jooken, Stephan Wagner}
\date{Source: arXiv:2508.20593v2 (CC BY 4.0)}
\begin{document}
\maketitle

\noindent\emph{Source and licence.} On the extrema of the mean subtree order of graphs. Version arXiv:2508.20593v2 (CC BY 4.0), distributed by arXiv under the Creative Commons Attribution 4.0 International licence, http://creativecommons.org/licenses/by/4.0/, as recorded in the arXiv licence metadata for this exact version and shown on the abstract page https://arxiv.org/abs/2508.20593v2. Full licence notice: \texttt{licenses/arxiv-2508.20593-CC-BY-4.0.txt}.\par\smallskip

\noindent\textit{Editorial scope.} Counted unit: the article's prop prop:localsubtreeorder (source0.tex lines 206--208). The article's complete original proof of it is delivered with it (source0.tex lines 210--234, one \texttt{proof} environment of 25 lines). No mathematics is written, repaired or re-derived: the statement and the proof are the article's own lines, in the article's own notation, and no retained line is substituted or re-flowed.  No further source-local context is needed: every cross-reference of every retained line resolves inside the two delivered spans.  Nothing was omitted from any retained span, so the package carries no omission record.  No retained line contains a citation, so the wrapper carries no bibliography and no bibliography entry.  No retained line contains a figure environment, an image-inclusion command or drawing code, so no image is delivered and the package declares no asset dependency.  Editorial formatting only, in addition: an AMS \texttt{amsart} preamble that restates the article's own declarations and macros used by the retained lines, namely the environments \texttt{prop} on separate counters, together with the standard packages those lines need.  The retained environments print in this document as Proposition 1 in order of appearance, whereas the article prints the same items with the numbering of its own full document.  The complete native source of the article as distributed in the arXiv e-print for this exact version is bundled unchanged as \texttt{sources/184030/source0.tex}.
\begin{prop}\label{prop:localsubtreeorder} If $T$ is a (non-empty) subtree of a connected graph $G$ of order $n$, then
\[\mu(G,T) \ge \frac{n+\abs{T}}{2}.\]
\end{prop}

\begin{proof} 
The proof is by induction on the integer $d = n- \abs{T}$. If $d=0$, then $\mu(G,T)=n=\frac{n+n}{2}$, so the base case of the induction is true. 
Assume that the statement is true for $d-1$, and let $(G,T)$ be a pair of a graph and a subtree for which $n-\abs T = d$. 
Let $Q= \bigcup_{v \in T} N(v) \setminus V(T)$ be the set of vertices that are adjacent to a vertex in $T$ but are not in $T$. 
We take a vertex $x$ in $Q$ and define $\mathcal T$ as the set of all minimal subtrees (with respect to inclusion) that contain both $T$ and $x.$ We let $y$ be a (fixed) neighbour of $x$ that belongs to $T$.
 

Let $G' = G \setminus x$, let $G_1$ be the connected component of $G \setminus x$ containing $T$, and let $G_2$ be the union of the other components. Then
\[N(G,T)= \sum_{S \in \mathcal T}  N(G, S )+ N(G',T) =\sum_{S \in \mathcal T} N(G, S ) + N(G_1,T)\]
and
\[R(G,T)= \sum_{S \in \mathcal T}  R(G, S )+ R(G',T) =\sum_{S \in \mathcal T} R(G, S ) + R(G_1,T).\]
Let us write $m$ for the number of vertices of $G_2$ %Denote the vertex set of $G_2$ by $V_2$, and let $m = \abs{V_2}$
(which is $0$ if $x$ is not a cut vertex). Observe that
$$\sum_{S \in \mathcal T} N(G, S ) \ge N(G, T \cup \{xy\}) \ge (m+1) N(G_1\cup \{xy\}, T \cup \{xy\}) = (m+1) N(G_1, T).$$
Here, the second inequality is true since there are at least $m+1$ ways to extend a subtree of $G_1\cup \{xy\}$ to a subtree of $G$.
For this, consider a spanning tree of $G \setminus G_1$ and a sequence of subtrees starting with $x$ and ending with the spanning tree, obtained by adding one edge at a time.


Now, using the induction hypothesis, we get
\begin{align*} R(G,T) &= \sum_{S \in \mathcal T} R(G, S ) + R(G_1,T)\\
&\ge \frac{n+(\abs{T}+1)}{2} \sum_{S \in \mathcal T} N(G, S ) + \frac{(n-1-m)+\abs T}{2} N(G_1,T) \\&=
\frac{n+\abs{T}}{2} \left( \sum_{S \in \mathcal T} N(G, S ) + N(G_1,T)  \right ) + \frac{1}{2} \Big (\sum_{S \in \mathcal T} N(G, S )- (m+1) N(G_1, T) \Big ) \\ & \ge \frac{n+\abs{T}}{2} N(G, T),
\end{align*}
completing the induction and thus the proof.
\end{proof} 

\end{document}
